\documentclass[../DoAn.tex]{subfiles}
\begin{document}

\section{Kết luận}

Đồ án đã xây dựng một hệ thống mô phỏng robot bao phủ trên bản đồ lưới có trọng số, có xét hướng di chuyển, chi phí quay, giới hạn năng lượng, cơ chế quay về căn cứ, sạc lại và vật cản động. Khác với bài toán tìm đường từ một điểm tới một điểm, bài toán bao phủ yêu cầu robot xử lý nhiều mục tiêu trong cùng một nhiệm vụ. Vì vậy, hệ thống không chỉ tìm đường tới ô chưa bao phủ tiếp theo, mà còn kiểm tra khả năng quay về căn cứ và trạng thái năng lượng sau mỗi quyết định.

Về mặt mô hình, môi trường được biểu diễn dưới dạng bản đồ lưới có trọng số, trong đó chi phí của mỗi ô là chi phí năng lượng tương đối đã chuẩn hóa. Robot được mô hình hóa bằng trạng thái có hướng $(v,h)$, thay vì chỉ xét tọa độ ô. Nhờ đó, thuật toán lập kế hoạch có thể phân biệt hai trường hợp robot đứng cùng một ô nhưng quay về hai hướng khác nhau và đưa chi phí quay vào quá trình tìm đường.

Về mặt thuật toán, hệ thống sử dụng Dijkstra có hướng để tìm đường chi phí nhỏ nhất trong không gian trạng thái mở rộng. Chi phí cạnh kết hợp chi phí đi vào ô kế tiếp và chi phí quay cần thiết để đi theo hướng mới. Trên nền tảng đó, robot chọn mục tiêu chưa bao phủ trong tập mục tiêu khả thi về năng lượng. Mỗi mục tiêu chỉ được chấp nhận nếu robot còn đủ năng lượng để đi tới mục tiêu và quay về căn cứ theo chính sách dự phòng của kịch bản.

Đồ án cũng xây dựng chính sách năng lượng quay về căn cứ có điều kiện. Với kịch bản không có vật cản động, hệ thống chỉ yêu cầu đủ năng lượng cho chặng đi tới mục tiêu và chặng quay về. Với kịch bản có vật cản động, hệ thống giữ thêm quỹ dự phòng $\max(M_{\min},M_{\text{prop}})$, trong đó thành phần tỷ lệ xấp xỉ $25\%$ chi phí quay về. Thiết kế này giúp chính sách năng lượng phản ánh đúng rủi ro của môi trường động, thay vì áp dụng cùng một mức dự phòng cho mọi bản đồ.

Phần thực nghiệm đã kiểm tra hệ thống trên bộ benchmark cuối cùng gồm 97 lần chạy, sau đó lọc còn 64 kịch bản duy nhất theo lần chạy mới nhất của từng \texttt{test\_name}. Kết quả cho thấy hệ thống hoàn thành tốt các bản đồ tĩnh liên thông, tăng số lần quay về khi năng lượng tối đa giảm, phân biệt đúng trạng thái \texttt{PARTIAL} khi bản đồ bị chia cắt, và ghi nhận \texttt{FAILED} khi vật cản động làm nhiệm vụ không thể kết thúc an toàn tại căn cứ. Các nhóm E và G cũng cho thấy chi phí địa hình và cấu trúc phân bố vật cản đều ảnh hưởng tới năng lượng và số bước, không chỉ riêng mật độ vật cản.

Từ các kết quả trên, có thể kết luận rằng đồ án đã đạt mục tiêu chính: xây dựng được một hệ thống mô phỏng robot bao phủ có xét chi phí năng lượng, chi phí quay, khả năng quay về căn cứ và phản ứng với vật cản động. Hệ thống không chỉ quan tâm tới việc đi qua nhiều ô nhất có thể, mà còn xét trạng thái nhiệm vụ sau cùng. Do đó, mô hình phù hợp với các nhiệm vụ có yêu cầu thu hồi hoặc sạc lại, nơi coverage cao nhưng không quay về được căn cứ không nên được xem là thành công đầy đủ.

Bên cạnh các kết quả đạt được, đồ án vẫn có một số giới hạn. Thứ nhất, hệ thống xét một robot duy nhất. Thứ hai, bản đồ được giả định là đã biết trong phạm vi mô phỏng. Thứ ba, mô hình năng lượng chưa xét đầy đủ động lực học, gia tốc, trượt bánh hoặc tiêu hao khi chờ. Thứ tư, vật cản động là các tác nhân không đối kháng. Thứ năm, chiến lược chọn mục tiêu là tham lam có ràng buộc, chưa tối ưu toàn cục thứ tự bao phủ.

\section{Vị trí của đồ án trong bối cảnh nghiên cứu và ứng dụng}

Đồ án không đặt đóng góp chính ở việc phát minh một thuật toán tìm đường hoàn toàn mới. Các thành phần nền tảng như Dijkstra, Coverage Path Planning, lập kế hoạch có ràng buộc năng lượng, cơ chế quay về căn cứ và lập kế hoạch lại trong môi trường thay đổi đều đã được nghiên cứu trong nhiều công trình trước đây. Đóng góp của đồ án nằm ở việc tích hợp các thành phần này thành một hệ thống mô phỏng hoàn chỉnh, có chương trình chạy được, có log, có dữ liệu thực nghiệm và có khả năng phân tích hành vi robot qua nhiều tình huống.

Trong hệ thống này, mô hình có chi phí quay được dùng làm mô hình chính, còn mô hình toàn hướng lý tưởng được dùng làm đối chứng. Cách tách này giúp đánh giá rõ ảnh hưởng của giả định cơ học tới hành vi lập kế hoạch. Khi chi phí quay được đưa trực tiếp vào trọng số cạnh, nó có thể làm thay đổi đường đi, số lần quay về, tổng năng lượng và trạng thái kết thúc nhiệm vụ. Vì vậy, nhóm H trong thực nghiệm không nhằm thay thế mô hình chính, mà giúp định lượng tác động của việc bỏ qua chi phí quay.

Từ góc nhìn ứng dụng, hệ thống phù hợp với hướng tự chủ có giám sát. Người vận hành không cần điều khiển từng bước di chuyển; robot tự xử lý các quyết định lặp lại như chọn mục tiêu, kiểm tra năng lượng, quay về căn cứ, sạc lại và lập kế hoạch lại khi môi trường thay đổi. Đây là nền tảng có thể mở rộng cho các nhiệm vụ kiểm tra, khảo sát hoặc vận hành trong khu vực có rủi ro.

\section{Các thử nghiệm phát triển không được giữ lại}

Trong quá trình phát triển, tác giả đã thử cài đặt một số heuristic tăng cường nhằm phát hiện sớm các vùng có nguy cơ bị bỏ sót sau lượt bao phủ đầu tiên. Ý tưởng chính là ưu tiên xử lý các ô hoặc cụm ô có dạng ``island'' trước khi robot rời khỏi khu vực lân cận, từ đó giảm số bước quay lại và giảm năng lượng tiêu thụ.

Tuy nhiên, việc xác định chính xác island không đơn giản. Island không nhất thiết là một ô đơn lẻ mà có thể là một cụm ô nhỏ, một đoạn hành lang hẹp, một hốc cục bộ hoặc một vùng phụ thuộc vào thứ tự bao phủ trước đó. Một số phiên bản heuristic cho hướng này đã làm chi phí tính toán tăng mạnh, trong khi mức cải thiện chưa ổn định trên các bản đồ khác nhau. Vì vậy, các heuristic này không được giữ lại trong phiên bản cuối. Quyết định này ưu tiên một hệ thống rõ ràng, ổn định và dễ kiểm chứng bằng benchmark.

\section{Hướng phát triển trong tương lai}

Từ các giới hạn trên, có thể mở rộng đồ án theo một số hướng. Hướng thứ nhất là xây dựng bộ chạy thử tự động ở chế độ không giao diện để đánh giá nhiều bản đồ và nhiều cấu hình hơn. Hướng thứ hai là cải tiến chiến lược chọn mục tiêu bằng các tiêu chí topo nhẹ, chẳng hạn mật độ vùng chưa bao phủ quanh mục tiêu hoặc chi phí quay về sau một cụm mục tiêu.

Hướng thứ ba là mở rộng mô hình năng lượng. Ngoài chi phí di chuyển và chi phí quay, có thể bổ sung chi phí chờ, chi phí tăng tốc, chi phí giảm tốc hoặc chi phí cảm biến. Hướng thứ tư là cải tiến mô hình vật cản động bằng dự đoán quỹ đạo ngắn hạn hoặc vùng an toàn theo bán kính. Hướng thứ năm là mở rộng sang nhiều căn cứ hoặc nhiều trạm sạc. Hướng thứ sáu là nghiên cứu phối hợp nhiều robot, trong đó các robot phải chia vùng bao phủ, tránh va chạm và chia sẻ trạng thái nhiệm vụ.

Tổng kết lại, đồ án đã xây dựng được nền tảng mô phỏng và phân tích cho bài toán bao phủ có ràng buộc năng lượng. Các hướng phát triển tiếp theo nên ưu tiên mở rộng từng thành phần một cách có kiểm soát và tiếp tục đánh giá bằng bộ benchmark rõ ràng.

\end{document}
